\begin{helper}
Fix \(n,m,k,S\), a real number \(E\), an edge parameter \(p\), and parameters
\(\varepsilon,\varepsilon_1,\varepsilon_2\).  Assume a deterministic witness
cover hypothesis: every sample satisfying the bad-on-regular event admits a
\(k\)-set \(I\) and a witness family \(B\) with the medium-class, mass,
cutoff, and normalized-coordinate lower-bound properties displayed below in
the proof.  Also assume that \(0<n\), that
every witness family \(B\) satisfying
\[
|B|\noderef{TwoBiteSmallCutoff}(\varepsilon,n)
\le
k n^{1/4-\varepsilon}+\noderef{TwoBiteLargeCutoff}(\varepsilon,n)
\]
has \(|B|\le S\), that the base-vertex universe has cardinality at most
\(n^2\), and that the count exponent \(E\) satisfies
\[
k\log n+2S\log n\le E.
\]
Then there is a finite candidate family indexed by a \(k\)-set
\(I\subseteq\mathrm{Fin}(n)\) together with an \(S\)-tuple of two-bite base
vertices.  Every sample in the bad-on-regular event is covered by one of these
candidates.  Moreover the number of candidate indices is at most
\[
\exp(E).
\]
\end{helper}
\begin{proof}
Let the index set consist of pairs \((I,f)\), where \(I\) is a finite subset
of \(\mathrm{Fin}(n)\) with \(|I|=k\), and
\[
f:\mathrm{Fin}(S)\to \noderef{TwoBiteBaseVertex}(m)
\]
is an \(S\)-tuple of base vertices.  The candidate associated to \((I,f)\)
is the event that there is a deterministic witness family \(B\) contained in
the image of \(f\), with the same medium-class, mass, cutoff, and normalized
coordinate lower-bound properties produced by
\(\noderef{MediumClosedPairsFailureDeterministicWitness}\).

For the cover, take a sample \(\omega\) in the bad-on-regular event.  The
assumed deterministic witness cover gives a witness \(I\) with \(|I|=k\), the
corresponding failed closed-pairs bound, and a nonempty witness family \(B\)
such that
\[
B\subseteq \noderef{TwoBiteMediumClass}(\omega,\varepsilon,I),
\]
\[
k n^{1/4-\varepsilon}
<
\sum_{x\in B}|\noderef{TwoBiteX}(\omega,I,x)|
\le
k n^{1/4-\varepsilon}+\noderef{TwoBiteLargeCutoff}(\varepsilon,n),
\]
\[
|B|\noderef{TwoBiteSmallCutoff}(\varepsilon,n)
\le
k n^{1/4-\varepsilon}+\noderef{TwoBiteLargeCutoff}(\varepsilon,n),
\]
and the normalized red/blue coordinate count is larger than
\[
\frac{k n^{1/4-\varepsilon}}{6(\log n)^2}.
\]
The size-conversion hypothesis gives \(|B|\le S\).  Since \(B\) is finite and
has at most \(S\) elements, choose an injection from \(B\) into
\(\mathrm{Fin}(S)\) and extend it to an \(S\)-tuple \(f\) whose image contains
\(B\).  Then \((I,f)\) is one of the candidate indices, and by construction
its candidate event holds for \(\omega\).

For the count, the number of choices for \(I\) is \(\binom nk\), which is at
most \(e^{k\log n}\) by \(\noderef{ParameterHierarchyT4ChooseExpBound}\).
The number of \(S\)-tuples of base vertices is
\(|\noderef{TwoBiteBaseVertex}(m)|^S\).  Since the base universe has size at
most \(n^2\), this is at most
\[
(n^2)^S=e^{2S\log n}.
\]
Multiplying these two estimates gives at most
\[
e^{k\log n}e^{2S\log n}=e^{k\log n+2S\log n}.
\]
Finally, the assumed inequality \(k\log n+2S\log n\le E\) and monotonicity of
the exponential give the claimed candidate-count bound.
\end{proof}
